\documentclass{article}
\usepackage{amsmath,amssymb}
\title{A restriction on three-letter clocks for deletion-two tag systems}
\author{}
\date{}
\begin{document}
\maketitle

\paragraph{Problem.}
A tagged block entered at alignment $p\in\{0,1\}$ produces the
concatenation of the appendants of its letters at positions $p,p+2,\ldots$
(zero-based). It passes alignment $p+|w|$ modulo two to the following
block. A one-sweep cell consists of two fixed finite words $w_0,w_1$
of even and odd length, respectively. An integrator outputs the word
$w_{s+p}$ from $w_s$, and a differentiator outputs $w_p$.
All bit arithmetic below is over $\mathbb F_2$ unless counts or ordinary
matrix eigenvalues are being computed.

\paragraph{Claim.}
A closed alphabet of three letters, one with empty appendant, cannot
implement both cells with the same productions. There is no bound on
appendant lengths or cell lengths in this statement.

\paragraph{Incidence reduction.}
Call the other letters $A,B$. Let $M$ be the two-by-two integer matrix
whose column $j$ counts the two letters in the appendant of $j$, and
let $c_s$ count them in $w_s$. For either cell the two alignments
produce the two different state words, once each. Consequently
\[
 Mc_0=Mc_1=c_0+c_1.
\]
If $M$ is invertible then $c_0=c_1$. The counts of the selected
letters in the two alignments must also agree since their images under
$M$ agree. Their outputs therefore have the same length, contradicting
the state parities. Thus $M$ is singular. Also $M(c_0+c_1)=2(c_0+c_1)$,
and $c_0+c_1\ne0$ since an eraser-only word cannot regenerate a cell.

\paragraph{Irreducible case.}
A nonnegative irreducible integer matrix of size two with eigenvalues
zero and two must have all entries one (use the trace and determinant).
Write $U,V$ for the two appendants. Each contains one $A$ and one $B$.
Every state word is an output and can therefore be parsed into copies
of $U,V$. The number $n$ of copies is the same in both states by the
incidence equation. The parsing used in subsequent sweeps can always
be the one induced by production, regardless of any ambiguity.

In a copy the two live letters lie in either the same lane or opposite
lanes. If both types use the same lane for their two letters, every
selected pair produces one $U$ and one $V$, fixing the output length
for fixed $n$. If just one type uses a single lane, exactly half the
copies of that type must be selected to preserve $n$ on each alignment.
Hence its multiplicity is even in either recurrent state. The other
type's multiplicity has parity $n\bmod 2$, fixing the total width
parity. Neither case permits opposite state parities.

It remains that both types have one live letter per lane. Their widths
must have opposite parities. Label each type by its width parity $s$.
It transmits $u+s$ on input $u$. The next type has bit $u+f(s)$
because the two alignments produce opposite types. Every Boolean
function of one bit is constant or $s$ plus a constant. On comparing
two input streams with the same initial cell state, the constant
drops out. Write $x$ for a one-sweep delay. The response of one copy
to an input difference is then either $1+x$ (constant $f$), or
$(1+x)^{-1}$ (nonconstant $f$). Every copy has the same response,
so a row of $n$ copies has its $n$-th power. Positive powers of one
choice cannot supply both $1+x$ and $(1+x)^{-1}$.

\paragraph{Reducible case.}
The diagonal entries are two and zero. Let $G$ denote the letter
with diagonal entry two. If its appendant $W$ contains no auxiliary
letter, no auxiliary letter can recur in a state word. Otherwise
triangularity forces the auxiliary letter to emit only erasers. Thus,
in the recurrent part, each sweep leaves copies of $W$ separated
by erasers. Each $W$ has two $G$s. The same number $n$ of copies
occurs in both states by the incidence equation for $G$ (after
discarding any nonrecurring auxiliary letter).

If both $G$s have the same parity within $W$, then precisely half
the copies must see each alignment to regenerate $n$ copies. This
fixes the total contribution from auxiliary letters too, since their
number selected is summed over equal numbers of each alignment.
The outputs have equal lengths.

Otherwise each $W$ regenerates one copy, surrounded by some erasers.
Assign the produced prefix and suffix to that same copy at each step.
Let the prefix and suffix parity functions of the alignment $p$ seen
by the core be $f(p)$ and $g(p)$. Consider differences between two
input histories and let $\alpha,\beta$ be the respective slopes of
these Boolean functions (zero or one). If $u$ is the incoming
difference before the prefix, $p$ the difference at the core, and $v$
the outgoing difference after the suffix, then
\[
 p=u+\alpha x p,\qquad v=p+\beta x p.
\]
The fixed width of $W$ cancels in differences. Every recurrent
copy consequently has the same transfer
$(1+\beta x)/(1+\alpha x)$. It is either $1$, $1+x$, or its inverse.
As before, powers of just one transfer cannot implement the two
required cells. The conclusion also covers zero auxiliary emissions.

\paragraph{Scope.}
This excludes a three-letter clock with a shared eraser under the
specified one-sweep cell interfaces. It places no restriction on
clocks operating over several sweeps, different clock interfaces,
or programs in which the clock and data use common live letters.
\end{document}
